% BHU PYQ -- Manually transcribed & error-corrected
% Paper: BCH-226 : Business Statistics
% Exam: B.Com. (Hons.) Semester IV Examination, 2022-23

\documentclass[11pt,a4paper]{article}
\usepackage[a4paper,top=2.5cm,bottom=2.5cm,left=2.8cm,right=2.8cm,headheight=14pt]{geometry}
\usepackage{amsmath,amssymb}
\usepackage[T1]{fontenc}
\usepackage[utf8]{inputenc}
\usepackage{lmodern,microtype,fancyhdr,enumitem,hyperref,lastpage,parskip}

\hypersetup{
    colorlinks=true,
    urlcolor=black,
    linkcolor=black,
    pdftitle={BCH-226: Business Statistics -- BHU B.Com. (Hons.) Sem IV 2022-23}
}

\pagestyle{fancy}
\fancyhf{}
\renewcommand{\headrulewidth}{0.4pt}
\renewcommand{\footrulewidth}{0.4pt}
\fancyhead[L]{\small\textit{Banaras Hindu University}}
\fancyhead[R]{\small\textit{BCH-226: Business Statistics}}
\fancyfoot[C]{\small Page \thepage\ of \pageref{LastPage}}

\newlist{parts}{enumerate}{2}
\setlist[parts,1]{label=(\alph*),leftmargin=*,itemsep=5pt,topsep=3pt}
\setlist[parts,2]{label=(\roman*),leftmargin=*,itemsep=3pt,topsep=2pt}
\newcommand{\pts}[1]{\hfill\textit{[#1]}}

\begin{document}

\begin{center}
  {\large\bfseries BANARAS HINDU UNIVERSITY}\\[2pt]
  {\normalsize B.Com. (Hons.) --- Semester IV Examination, 2022-23}\\[6pt]
  \rule{0.8\linewidth}{0.5pt}\\[6pt]
  {\large\bfseries Paper: BCH-226: Business Statistics}\\[4pt]
  {\normalsize Time: 3 Hours\hfill Maximum Marks: 70}
\end{center}

\rule{\linewidth}{0.4pt}

\smallskip
\noindent\textit{Instructions: Attempt all questions. All questions carry equal marks (14 marks each).}

\bigskip

\noindent\textbf{Question 1.} \\
``Index numbers are a specialised type of averages.'' Comment on this statement. Also state briefly the various points to be considered in the construction of index numbers.\pts{14}

\centerline{\textbf{OR}}

The following are the average daily wages and price indices of a group of industrial workers:
\begin{center}
\renewcommand{\arraystretch}{1.2}
\small
\begin{tabular}{|c|c|c|}
\hline
\textbf{Year} & \textbf{Average Daily Wages (Rs.)} & \textbf{Price Index} \\ \hline
1999          & 80                                 & 100                  \\
2000          & 108                                & 120                  \\
2001          & 125                                & 125                  \\
2002          & 147                                & 140                  \\
2003          & 216                                & 180                  \\
2004          & 230                                & 200                  \\ \hline
\end{tabular}
\end{center}
\begin{parts}
  \item Compute the index number of real wages.\pts{8}
  \item In which year do the workers have maximum buying power?\pts{6}
\end{parts}

\medskip\rule{\linewidth}{0.2pt}

\noindent\textbf{Question 2.} \\
What do you mean by ``Cost of Living Index''? State its kinds. Explain any one of its kinds with the help of imaginary figures.\pts{14}

\centerline{\textbf{OR}}

Compute Fisher's Ideal Index Number from the following data:
\begin{center}
\renewcommand{\arraystretch}{1.2}
\small
\begin{tabular}{|c|c|c|c|c|}
\hline
\textbf{Items} & \multicolumn{2}{c|}{\textbf{Base Year}} & \multicolumn{2}{c|}{\textbf{Current Year}} \\ \cline{2-5}
               & \textbf{Price} & \textbf{Quantity}      & \textbf{Price} & \textbf{Quantity}         \\ \hline
A              & 10             & 40                     & 12             & 50                        \\
B              & 12             & 25                     & 15             & 20                        \\
C              & 15             & 10                     & 20             & 12                        \\
D              & 20             & 5                      & 30             & 2                         \\ \hline
\end{tabular}
\end{center}
Also examine if it satisfies the Time Reversal Test and Factor Reversal Test.\pts{14}

\medskip\rule{\linewidth}{0.2pt}

\noindent\textbf{Question 3.} \\
What are the components of a Time Series? Explain the significance of moving averages in analyzing time series and point out its limitations.\pts{14}

\centerline{\textbf{OR}}

The following data relates to the sales of Dinkar Limited:
\begin{center}
\renewcommand{\arraystretch}{1.2}
\small
\begin{tabular}{|l|c|c|c|c|c|}
\hline
\textbf{Year}                 & 2018 & 2019 & 2020 & 2021 & 2022 \\ \hline
\textbf{Sales (Rs. in lakhs)} & 80   & 120  & 160  & 360  & 280  \\ \hline
\end{tabular}
\end{center}
\begin{parts}
  \item Fit a straight line trend by the method of least squares and tabulate the trend values.\pts{6}
  \item Estimate the sales for the year 2023.\pts{4}
  \item Measure the short-term fluctuations.\pts{4}
\end{parts}

\medskip\rule{\linewidth}{0.2pt}

\noindent\textbf{Question 4.} \\
What are the various methods of measuring seasonal variation in a time series? Discuss any one method in detail.\pts{14}

\centerline{\textbf{OR}}

Calculate the seasonal index by the link relative method from the following link relatives:
\begin{center}
\renewcommand{\arraystretch}{1.2}
\small
\begin{tabular}{|c|c|c|c|c|c|}
\hline
\textbf{Quarter} & \textbf{2018} & \textbf{2019} & \textbf{2020} & \textbf{2021} & \textbf{2022} \\ \hline
I                & --            & 80            & 88            & 80            & 83            \\
II               & 120           & 117           & 129           & 125           & 117           \\
III              & 133           & 113           & 111           & 115           & 120           \\
IV               & 83            & 89            & 93            & 96            & 79            \\ \hline
\end{tabular}
\end{center}\pts{14}

\medskip\rule{\linewidth}{0.2pt}

\noindent\textbf{Question 5.} \\
What do you understand by ``Statistical Quality Control''? Discuss its techniques and advantages.\pts{14}

\centerline{\textbf{OR}}

The mean and range of 10 random samples of size 4 are as follows:
\begin{center}
\renewcommand{\arraystretch}{1.2}
\small
\begin{tabular}{|l|c|c|c|c|c|c|c|c|c|c|}
\hline
\textbf{Sample No.}           & \textbf{1} & \textbf{2} & \textbf{3} & \textbf{4} & \textbf{5} & \textbf{6} & \textbf{7} & \textbf{8} & \textbf{9} & \textbf{10} \\ \hline
\textbf{Sample Mean ($\bar{X}$)} & 494        & 508        & 500        & 582        & 552        & 538        & 514        & 614        & 600        & 558         \\ \hline
\textbf{Sample Range ($R$)}    & 95         & 128        & 100        & 91         & 68         & 65         & 148        & 28         & 37         & 80          \\ \hline
\end{tabular}
\end{center}
Determine the control limits for the $\bar{X}$ chart. Draw the three lines of the mean chart on graph paper. Also determine whether the process is within control or not? You are given that $d_2 = 2.059$ for a sample size of 4.\pts{14}

\vfill
\rule{\linewidth}{0.4pt}
\begin{center}\small
  \href{https://bhu-exam.vercel.app/}{Click here to visit our website}\\[4pt]
  \href{https://me-aryan.vercel.app/}{Click here to visit our contributor}\\[4pt]
  \href{https://quashan-me.vercel.app/}{Click here to visit our contributor}
\end{center}

\end{document}
